\section{Lý thuyết trò chơi}
\begingroup\small
\begin{multicols}{2}
\subsection{Trạng thái thắng, thua và hòa}
Xét trò chơi hai người, thông tin hoàn hảo, không ngẫu nhiên, cùng
tập nước đi cho cả hai người; ai \emph{không đi được} ở lượt mình
thì thua (quy tắc chơi chuẩn). Với đồ thị trạng thái không chu trình,
đặt $W(s)$ là trạng thái thắng của người sắp đi. Trạng thái không
có nước đi là thua; trạng thái có ít nhất một cạnh tới trạng thái
thua là thắng; nếu mọi cạnh tới trạng thái thắng thì là thua.
Tính theo thứ tự topo ngược. Nếu cần giá trị điểm thay vì thắng/thua,
dùng minimax với quy ước giá trị nhất quán theo người chơi.

\subsection{Nim}
Với các đống có kích thước $a_1,\ldots,a_k$, mỗi lượt giảm đúng
một đống đi một lượng dương. Gọi $x=a_1\oplus\cdots\oplus a_k$.
Trạng thái thua khi và chỉ khi $x=0$. Nếu $x=0$, mọi nước đi
đều làm XOR khác $0$. Nếu $x\ne0$, chọn bit cao nhất bật trong
$x$, tìm đống $a_i$ có bit ấy bật, và đổi nó thành
$a_i'=a_i\oplus x<a_i$ để XOR mới bằng $0$.
Kết luận này dùng quy tắc chơi chuẩn. Với misère Nim (ai lấy viên
cuối cùng thì thua): nếu mọi đống khác rỗng đều có kích thước $1$,
người sắp đi thua khi số đống đó lẻ; nếu còn đống lớn hơn $1$,
tiêu chuẩn thắng vẫn là XOR khác $0$.

\subsection{Số Grundy và mex}
Với trò chơi vô tư hữu hạn, không chu trình, định nghĩa
\[
g(s)=\operatorname{mex}\{g(t):s\to t\},\qquad
\operatorname{mex}(S)=\min\{z\ge0:z\notin S\}.
\]
Trạng thái hết nước đi có $g=0$; $g(s)=0$ đúng khi $s$ thua.
Nếu một lượt chỉ được đi trong \emph{một} trong nhiều trò chơi
độc lập, định lý Sprague--Grundy cho
$g_{\rm tổng}=g_1\oplus\cdots\oplus g_k$.
Không dùng XOR này khi một nước đi có thể thay đổi đồng thời nhiều
thành phần, hoặc khi luật chơi là misère mà chưa biến đổi riêng.

\subsection{Trò chơi trên đồ thị có chu trình}
Trên đồ thị hữu hạn có chu trình, quy hoạch động topo không còn
hợp lệ. Đánh dấu các đỉnh hết nước là thua, rồi lan ngược:
tiền nhiệm của đỉnh thua là thắng; tiền nhiệm mà \emph{mọi} nước đi
đều tới đỉnh thắng là thua. Các đỉnh chưa được đánh dấu có thể
duy trì ván đấu mãi mãi, nên là hòa nếu chơi tối ưu theo quy ước
không kết thúc là hòa. Cần xác định rõ luật của đề: trò chơi có
giới hạn lượt hoặc có trạng thái lặp bị xử thua sẽ cho kết quả khác.

\paragraph{Mẫu C++17.}
\texttt{grundyDAG} nhận thứ tự topo; \texttt{gameOutcome} trả
$0$ thua, $1$ thắng, $2$ hòa. Hai hàm nhận đỉnh $[0,n)$.
\cppsnippet{code/game_theory/game.hpp}
\end{multicols}
\endgroup
